\documentclass[11pt]{article}

\usepackage{amsmath,amssymb}
\usepackage{xurl}
\usepackage[hidelinks]{hyperref}
\setlength{\emergencystretch}{2em}

% Shared commands for notes in docs/paper-gaps/.

\DeclareUnicodeCharacter{2080}{\ifmmode{}_{0}\else\textsubscript{0}\fi}
\DeclareUnicodeCharacter{2081}{\ifmmode{}_{1}\else\textsubscript{1}\fi}
\DeclareUnicodeCharacter{2082}{\ifmmode{}_{2}\else\textsubscript{2}\fi}
\DeclareUnicodeCharacter{2099}{\ifmmode{}_{n}\else\textsubscript{n}\fi}

\newcommand{\Lean}{\textsc{Lean}}
\ExplSyntaxOn
\NewDocumentCommand{\leanid}{m}{
  \ifmmode
    \textsf{\detokenize{#1}}
  \else
    \group_begin:
    \urlstyle{sf}
    \tl_set:Nn \l_tmpa_tl {#1}
    \tl_replace_all:Nnn \l_tmpa_tl {\_} {_}
    \exp_args:NV \path \l_tmpa_tl
    \group_end:
  \fi
}
\ExplSyntaxOff
\providecommand{\tr}{\operatorname{tr}}
\newcommand{\tnleanIssueBase}{https://github.com/LionSR/TNLean/issues/}
\newcommand{\tnleanPrBase}{https://github.com/LionSR/TNLean/pull/}
\newcommand{\tnleanDocBase}{https://sirui-lu.com/TNLean/docs/}
\newcommand{\ghissue}[1]{\href{\tnleanIssueBase#1}{GitHub issue \##1}}
\newcommand{\ghpr}[1]{\href{\tnleanPrBase#1}{PR \##1}}
\newcommand{\doclink}[1]{\href{\tnleanDocBase#1.html}{\path{#1}}}

% Dirac notation and the identity operator, as in blueprint/src/macros/common.tex.
% \providecommand keeps note-local definitions authoritative.
\usepackage{dsfont}
\providecommand{\ket}[1]{|#1\rangle}
\providecommand{\bra}[1]{\langle #1|}
\providecommand{\braket}[2]{\langle #1 | #2 \rangle}
\providecommand{\ketbra}[2]{|#1\rangle\!\langle #2|}
\providecommand{\Id}{\mathds{1}}
\providecommand{\one}{\mathds{1}}
\providecommand{\C}{\mathbb{C}}
\providecommand{\MN}[1]{M_{#1}(\C)}
\providecommand{\MD}{M_D(\C)}

% External referencing for paper sources.  The build helper writes sanitized
% auxiliary files under build/paper-aux/ and excludes duplicated source labels.
\usepackage{xr}

% Import a generated paper auxiliary file with a caller-supplied label prefix.
% The candidate paths support builds from docs/paper-gaps/, the repository root,
% and build/paper-aux/ itself.
\newcommand{\importpaperaux}[2]{%
  \IfFileExists{../../build/paper-aux/#2.aux}{%
    \externaldocument[#1]{../../build/paper-aux/#2}%
  }{%
    \IfFileExists{build/paper-aux/#2.aux}{%
      \externaldocument[#1]{build/paper-aux/#2}%
    }{%
      \IfFileExists{#2.aux}{%
        \externaldocument[#1]{#2}%
      }{}%
    }%
  }%
}

% General external reference macros
\newcommand{\paperref}[2]{\ref{#1-#2}}
\newcommand{\papereqref}[2]{(\ref{#1-#2})}

% CPSV16 source (1606.00608 / MPDO-22-12-17-2.tex) external document and shortcuts
\importpaperaux{cpsv16-}{1606.00608/MPDO-22-12-17-2-tnlean}
\newcommand{\cpsvref}[1]{\paperref{cpsv16}{#1}}
\newcommand{\cpsveqref}[1]{\papereqref{cpsv16}{#1}}

% Verdict marker: \gapnote{<kind>}{<status>}. Kind is the policy
% classification (clarification, local-correction, scope-restriction,
% unfaithful, false-source, open-gap); status is open, wip (elimination
% underway), resolved, or historical. Machine-read by the site index;
% prints a verdict line.
\newcommand{\gapnote}[2]{\par\noindent\textbf{Verdict:} #1 (#2).\par}


\title{Paper Comparison for BNT Inputs in the Non-periodic MPS Fundamental Theorem}
\author{}
\date{2026-05-05}

\begin{document}
\maketitle

\section{The source comparison}

The non-periodic matrix product vector Fundamental Theorem in
\cite{Cirac2016MPDO_arXiv,Cirac2017MPDO_AnnPhys,Cirac2021Matrix} is a theorem
about tensors already put in canonical form. The comparison is not stated as a
theorem about arbitrary block-diagonal data plus unrelated finite-length
hypotheses. The finite-length hypotheses now isolated in the formalization are
not separate assumptions in the paper theorem. They are placeholders for
specific source arguments: canonical-form reduction, BNT selection, fixed-length
block-injective span after blocking, phase-gauge matching, and power-sum
comparison.

The source matrix product vectors are indexed by positive lengths.  The 2016
paper defines \( |V^{(N)}(A)\rangle \) for \(N=1,2,\ldots\) and compares
families at those lengths \cite[lines~130--151]{Cirac2016MPDO_arXiv}.  Thus a
comparison relation that also includes the empty word is slightly stronger than
the relation appearing in the paper.  The length-zero term is an auxiliary
convention that must be handled separately when zero blocks or cumulative spans
are present.

The first source step removes invariant subspaces. Starting from a tensor
\(B\), the 2016 paper replaces it by block-diagonal matrices
\[
  A^i=\bigoplus_{k=1}^r \mu_k A_k^i
\]
which generate the same positive-length matrix product vectors
\cite[lines~214--225]{Cirac2016MPDO_arXiv}. Each block is then treated through
its completely positive map. Periodic peripheral eigenvalues are removed by
blocking a common multiple of their periods
\cite[lines~227--231]{Cirac2016MPDO_arXiv}. After that blocking, the blocks are
normal in the sense used in the paper, and the tensor is in canonical form
\cite[lines~233--255]{Cirac2016MPDO_arXiv}. The review gives the same
description: block diagonalization, period-removing blocking, primitive
transfer maps, and the resulting canonical form
\cite[lines~1793--1839]{Cirac2021Matrix}.

The next source step chooses a basis of normal tensors. In the 2016 paper, such
a basis \(A_1,\ldots,A_g\) is a family of normal tensors whose matrix product
vectors span the matrix product vectors of \(A\) at each length and are
eventually linearly independent
\cite[lines~271--274]{Cirac2016MPDO_arXiv}. The subsequent characterization
identifies the basis tensors with representatives of the phase-gauge
equivalence classes of normal blocks: every normal block in the canonical form
is a phase times an invertible conjugate of one selected basis tensor, and no
two selected basis tensors are related in that way
\cite[lines~278--280]{Cirac2016MPDO_arXiv}. The review repeats the same
definition and characterization
\cite[lines~1846--1859]{Cirac2021Matrix}.

For a tensor in canonical form, the source then writes the matrix product vector
family through the basis of normal tensors as
\[
  |V^{(N)}(A)\rangle
    = \sum_{j=1}^g
        \left(\sum_{q=1}^{r_j}\mu_{j,q}^N\right)
        |V^{(N)}(A_j)\rangle.
\]
These are the data compared between two tensors in the equal or proportional
case
\cite[lines~283--302]{Cirac2016MPDO_arXiv};
\cite[lines~1864--1885]{Cirac2021Matrix}.

The Fundamental Theorem itself says that, for two tensors \(A\) and \(B\) in
canonical form, proportionality of all matrix product vector families matches
their bases of normal tensors up to a permutation, phases, and invertible gauge
matrices
\cite[lines~347--352]{Cirac2016MPDO_arXiv};
\cite[lines~1890--1900]{Cirac2021Matrix}. In the equal case, the source then
compares the coefficients in the BNT decompositions by the power-sum identity
\[
  \sum_{q=1}^{r_{a,j}}\mu_{j,q}^N
  =
  \sum_{q=1}^{r_{b,j}}(\nu_{j,q}e^{i\phi_j})^N
\]
for each matched basis tensor and all lengths \(N\). The elementary power-sum
lemma recovers the multiplicities and weights, and the blockwise gauges assemble
into the global conjugating matrix
\cite[lines~1155--1192]{Cirac2016MPDO_arXiv}.

\section{Fixed-length spans after blocking}

There are two different blocking operations in the source account. The
period-removing blocking makes the transfer maps primitive. The Wielandt
blocking supplies finite-length injectivity after normality is available. These
steps should not be merged.

For one normal tensor, the quantum Wielandt theorem is cited as saying that
after at most \(D^4\) blockings the tensor becomes injective
\cite[line~332]{Cirac2016MPDO_arXiv}. The older canonical-form paper defines
injectivity by the map
\[
  \Gamma_L(X)
    =
    \sum_{i_1,\ldots,i_L}
      \operatorname{tr}(X A_{i_1}\cdots A_{i_L})
      |i_1,\ldots,i_L\rangle,
\]
and observes that injectivity of \(\Gamma_L\) is equivalent to the length-\(L\)
word products spanning the whole matrix algebra
\cite[lines~888--908]{PerezGarcia2007MPSReps}. Thus the relevant conclusion is
a homogeneous fixed-length span statement after a positive blocking length, not
mere membership in a cumulative span. The 2021 review records a sharper
MPS-specific Wielandt bound for normal tensors
\cite[lines~1942--1945]{Cirac2021Matrix}; the point here is the nature of the
conclusion, not which numerical bound is ultimately used.

For several basis blocks, the paper needs block-injective canonical form. It
defines this as the statement that one common physical tensor spans every block
matrix simultaneously:
\[
  \bigl(A_1^i,\ldots,A_g^i\bigr)_i
  \quad\text{spans}\quad
  \bigoplus_{j=1}^g M_{D_j}(\mathbb C)
\]
after the appropriate blocking
\cite[lines~317--332]{Cirac2016MPDO_arXiv}. The source then cites the
P\'erez-Garc\'ia--Verstraete--Wolf--Cirac bound: if \(L\) is a length at which
each normal tensor is injective, then after \(3(L+1)(D-1)\) blockings the
canonical-form tensor is block-injective; together with \(L\le D^4\), this gives
the \(3D^5\) bound
\cite[lines~340--345]{Cirac2016MPDO_arXiv}. The review states the same
block-injective conclusion and bound
\cite[lines~2114--2124]{Cirac2021Matrix}.

This block-injective conclusion is stronger than separate injectivity of the
individual blocks. It includes the cross-block separation needed to prescribe
different matrices in different basis components at one common length. In dual
form, one fixed word length \(L\), after the prescribed blocking, gives the
following implication. If, for every word \(w\) with \(|w|=L\),
\[
  \sum_{j=1}^g \operatorname{tr}(\Delta_j A_j^w)=0
\]
then every test matrix \(\Delta_j\) is zero. Equivalently, the homogeneous
length-\(L\) word evaluations span the direct sum
\[
  \bigoplus_{j=1}^g M_{D_j}(\mathbb C).
\]
This fixed-length trace-dual statement is the mathematical content behind the
pair trace-separation predicates used in the non-periodic proof. It is not a
separate theorem in the papers; it is the dual form of block-injective
direct-sum span. In the later MPDO argument the same source explicitly writes
\(C_L(V)=\operatorname{span}(V^L)\) and invokes the block-injective proposition
to obtain \(C_L(V)=\mathcal A_1\) for some fixed \(L\)
\cite[lines~1984--1988]{Cirac2016MPDO_arXiv}. This is again an exact-length
span statement.

\section{The homogeneous padding point}

The paper-level injectivity and block-injectivity conclusions are homogeneous:
they concern words of a specified positive length after blocking. They are not
consequences of the empty word lying in a cumulative word span.

This distinction matters for the pair representation used in the comparison
proof. A cumulative span includes the empty word. Therefore the identity pair
belongs to a cumulative span at length zero for every pair of tensors. That fact
gives no positive-length padding. For example, over a one-letter alphabet with
\(A_0=B_0=0\) in bond dimension one, the length-one pair span is generated by
the zero pair, while the identity pair appears at length zero. Hence the
implication from cumulative identity membership to identity membership in all
later homogeneous spans is false.

Cumulative fullness alone is still too weak. In bond dimension one with one
letter and pair word values \((2^n,3^n)\), the cumulative span through length
one is all of \(\mathbb C\oplus\mathbb C\), since \((1,1)\) and \((2,3)\) are
linearly independent. At each positive homogeneous length \(n\), however, the
span is only the line through \((2^n,3^n)\), so it does not contain \((1,1)\).
Thus the source hypotheses must not be discarded down to cumulative pair-span
fullness.

The source target is fixed-length block-injective direct-sum span after the
prescribed blocking. A Burnside--Jacobson homogeneous-padding argument would be
an alternative proof of the same kind of exact-length trace separation; it must
retain the BNT, non-gauge-equivalence, and normalization data that make the
fixed-length conclusion true. It is not a standalone consequence of full
cumulative pair-span, and it is not a restatement of the empty-word observation.

\section{Comparison data after common blocking}

The equal-case source proof compares BNT decompositions after the canonical form
has been chosen. The after-blocking theorem reaches this situation through a
longer route: remove the zero tail, choose one common physical blocking length
for both tensors, reindex nested blocked words against flat words, separate
period removal from later injectivity blocking, and then compare the resulting
nonzero sector families.

The common-blocking issue is mathematical rather than clerical. If a block has
period-removal length \(m_k\) and the final common length is \(p=m_ke_k\), the
tensor obtained by blocking first by \(m_k\) and then by \(e_k\) has nested
physical labels, whereas the tensor obtained by blocking once by \(p\) has flat
word labels. Paper notation treats these identifications silently; the
theorem must prove the reindexing statements.

The BNT comparison data should therefore be produced for the final nonzero
sector families after all chosen blockings and reindexings. In source terms,
this means the proof must construct the matched basis tensors, the phase-gauge
equivalences, and the power-sum comparison for exactly those families. The
common primitive sectors should first be grouped into representative BNT
classes, with repeated cyclic representatives and transported weights accounted
for by the power-sum comparison. Common phase covers and proportional
decomposition data match the source theorem only for those representative or
grouped BNT families, not for the raw flattened cyclic-sector list. The raw list
can contain several cyclic representatives of the same original block with the
same transported weight, so it is not the basis of normal tensors used in the
paper comparison.

\section{Current mathematical obligations}

The current formalization has conditional theorems at the following mathematical
points.

\begin{itemize}
  \item The homogeneous pair identity-padding route must be justified by a
  source-faithful fixed-length argument. A theorem may assume a period and
  residue window as an intermediate statement, but the missing mathematical
  content is to obtain that window, or an equivalent fixed-length direct-sum
  span, from the BNT/non-gauge/normalization situation.

  \item The unconditional BNT trace-separation theorem is the dual form of the
  fixed-length direct-sum span supplied by block-injective canonical form. It
  remains open until this direct-sum span is proved for every ordered pair of
  distinct basis blocks and one common length is chosen by a finite maximum.

  \item Common phase-cover and proportional-decomposition data must be obtained
  for the final common nonzero sector families from BNT matching and the
  power-sum comparison. Passing those data as unrelated assumptions does not
  prove the paper theorem.

  \item The after-blocking comparison needs the paper-facing inputs for the
  representative BNT families: positive dimensions, injectivity, normalization,
  self-overlap limits, off-overlap decay for inequivalent representatives, and
  equality of finite-length matrix product vector spans. The auxiliary structure
  used to collect these inputs is only a formal device; the mathematical
  obligation is to derive the listed assertions for the representatives selected
  by the canonical-form and BNT construction.
\end{itemize}

For traceability, these obligations are being followed in \ghissue{1133},
\ghissue{1178}, \ghissue{1068}, and \ghissue{944}, with \ghissue{1170} tracking
the larger non-periodic proof. The issue numbers are not part of the
mathematical statement.

\section{Closing criterion}

A proof issue in this part of the theorem should be closed only when the
following three assertions agree.

First, the source assertion has been identified at the level of the mathematical
objects above: canonical form, basis of normal tensors, block-injective
fixed-length span, phase-gauge BNT matching, or power-sum comparison. Second,
the formal theorem states the same assertion for the same blocked and reindexed
families, or explicitly records the extra hypothesis that remains. Third, the
final comparison theorem uses that assertion rather than an unrelated assumption
with a similar shape.

Conditional theorems are useful because they expose the missing mathematical
inputs precisely. They should not be described as closing the source theorem
until the corresponding assertion has been proved or the final theorem is
deliberately restated with that hypothesis.

\bibliographystyle{alpha}
\bibliography{docs/paper-gaps/references}

\end{document}
